\section{Einleitung}
\label{sec:Einleitung}
Laut Definition ist \glqq Eine Datenbank ist eine Sammlung von Daten, die untereinander in einer logischen Beziehung stehen und von einem eigenen Datenbankmanagementsystem (Database Management System, DBMS) verwaltet werden\grqq (\cite{kaufmannSQLNoSQLDatenbanken92023}, S.3). Das DBMS ermöglicht idealerweise einen Zugriff durch mehrere Anwendungen oder Endnutzer, welche gleichzeitig, fehlerfrei, schnell und sicher die Informationen aus den Daten abfragen und konfliktfrei verändern können, und das bei minimalem Bedarf an Systemressourcen.

Diese Vorteile einer zentralen Datenverwaltung für mehrere Klienten werden insbesondere deutlich, wenn man das Datenbanksystem mit einer einfachen Dateiverwaltung vergleicht. Hier verwendet jede Anwendung und Endnutzer ein eigenes Dateisystem und physische Ablage auf einem Speichersystem. Redundanzen und Inkonsistenzen entstehen schnell, wenn die gleichen Daten durch Kopieren an einen anderen Speicherplatz übertragen oder redundante Daten nur an einer Stelle verändert werden. Auch ist die logische Struktur der Dateiablage und Codierungen (Dateiformate) individuell auf die Anwendungen abgestimmt, was die Flexibilität bei einer erforderlichen gemeinsamen Datenverarbeitung beeinträchtigt.

Je nach Anwendung werden unterschiedliche Datenbanksysteme eingesetzt, die sich deutlich in der Datenstruktur unterscheiden. In der heutigen Zeit haben sich insbesondere das relationale Modell und No-SQL-Datenbanken (im Bereich von Big-Data) durchgesetzt, während hierarchische Modelle und objektorientierte Modelle (oder Mischformen wie objektrelationale Modelle) weitestgehend verdrängt wurden (\cite{kaufmannSQLNoSQLDatenbanken92023}: 12ff; \cite{gehringWirtschaftsinformatik2022}: 781). Das relationale Modell bietet im Vergleich zu den anderen Modellen insbesondere den Vorteil, dass die Datenkonsistenz sehr gut eingehalten werden kann, da die Datentypen streng einheitlich definiert werden.

Das Vorgehen bei der Erstellung einer Datenbank ist die \textit{Datenbankmodellierung} und verfolgt die Ziele (\cite{schickerDatenbankenUndSQL2025}, S.23) einer
\begin{itemize}[itemsep=3pt]
	%https://www.overleaf.com/learn/latex/Lists#The_itemize_environment_for_bulleted_(unordered)_lists
	\item{geringen Redundanz,} 
	\item{guten Bedienbarkeit und einfachen Erweiterbarkeit,} 
	\item{schnellen Zugriffszeiten,} 
	\item{Konsistenz,}
	\item{und Integrität.}
\end{itemize}

Die Ziele können teils miteinander konkurrieren, sodass es oft nicht möglich ist, alle Ziele gleichsam zu optimieren. Beispielsweise ist vorstellbar, dass aus dem Datenbestand abgeleitete oder berechnete Daten für eine bessere Performance in verarbeiteten Datentypen zusätzlich gespeichert werden. Die eindeutige funktionale Beziehung zu den Rohdaten erzeugt somit eine Redundanz, welche außerdem die Datenkonsistenz im Fehlerfall beeinträchtigen kann. Weiterhin ist es möglich, dass im Sinne einer Redundanz-Optimierung gewisse Attribute, welche in verschiedenen Objekten der Datenbank in identischer Form vorkommen, separat gespeichert werden. Werden nun die Objekte gelöscht, bleiben diese separat gespeicherten Attribute womöglich bestehen, wodurch die referentielle Integrität beeinträchtigt wird, da Datenbestände entstehen, die keinen sinnvollen Bezug mehr zu den eigentlich für die Datenbank wichtigen Datenbeständen haben. Diese Konflikte zeigen die Schwierigkeit bei den Zielstellungen einer Datenbankmodellierung. Es ist daher wichtig, dass die Datenbankmodellierung nicht nur an den bestehenden Bestand von Datensätzen angepasst ist, sondern auch wachsende Anforderungen an die Anzahl der Datensätze sowie zusätzliche Funktionen von Anfang an mitgedacht werden. Die flexible Erweiterbarkeit sollte daher bei potenziell wachsenden Datenbanken stets mitgedacht werden.

\subsection{Problemstellung}
\label{subsection:Problemstellung}
In dieser Arbeit wird die Datenbankmodellierung an einem selbstgewählten Anwendungsfall schrittweise durchgeführt. Beispielhaft wird die Umsetzung eines Datenbankmodells anhand der Carsharing App \glqq TeilAuto Neckaralb\grqq demonstriert, die vom Autor dieser Arbeit regelmä{\ss}ig verwendet wird. Ähnlich anderen Carsharing-Anbietern ermöglicht diese App die zeitlich und örtlich flexible Miete diverser Fahrzeuge und Fahrzeugklassen vom Lastenrad, über das Familienauto bis hin zum Umzugstransporter. Monatlich wird für jede/n Kund:in eine Rechnung erstellt, die die persönlichen Buchungen zusammen mit dem im Vertrag festgelegten Tarif verrechnet. Hierbei können mehrere Kund:innen den gleichen Vertrag besitzen wenn sie diesem als Gruppe (Haushalt, Familie) zugeteilt sind und so einen vergünstigten Gruppen-Tarif erhalten, wobei die Rechnungen weiterhin kund:innenspezifisch erstellt werden. Das Datenmodell soll sowohl die Sichtweise von Kund:innen, als auch dem Carsharing-Anbieter widerspiegeln und es ermöglichen, dass Kund:innen bei Buchungen auf die für sie relevanten Informationen zugreifen können, als auch die Mitarbeitenden des Unternehmens die Informationen erhalten um Kund:innen, Verträge, Abrechnungen und Fahrzeuge zu managen.

\subsection{Ziel und Aufbau der Arbeit}
Das Hauptziel dieser Arbeit ist es, den vollständigen Entwurfsprozess einer Datenbank anhand des Beispiels der Carsharing-Applikation zu erklären und praktisch zu erarbeiten. Zuerst werden in \autoref{section:basics} die theoretischen Grundlagen zur Modellierung relationaler Datenbanken erläutert. Als grundlegende Schritte finden sich hier die Anforderungsanalyse, aus der ein Entity-Relationship-Modell (ERM) konzipiert wird, das die grundlegenden Entitäten, Attribute und Beziehungen grafisch darstellt und die Überführung des ERM in ein in dieser Arbeit gewähltes Relationenmodell, das die Entitätstypen und Beziehungstypen in Tabellen überführt. Es wird gezeigt, durch welche Normalisierungs-Regeln diese Tabellen bezüglich der Ziele einer Datenbankmodellierung optimiert werden können. Anschließend wird im praktischen Teil in \autoref{section:main} das Beispiel der Carsharing-Applikation auf Basis der vorgestellten Modellierungsschritte vollständig entworfen und in der Programmiersprache Racket und dem DBMS SQLite praktisch umgesetzt. Abschließend wird in \autoref{section:summary} die Arbeit zusammengefasst und die Umsetzung kritisch beleuchtet.
